The failures concern different mathematical quantities. A proposal can put
particles in the wrong place, an importance correction can leave highly
variable weights, and an accurately evaluated derivative can still estimate
the model score poorly. Treating these as one problem obscures which repair
an experiment has actually tested.

\subsection{A correct flow applied to the wrong conditional distribution}

A UKF describes predictive uncertainty using moments. An LEDH map then uses
the observation to move particles. Neither fact guarantees that the map's
reference covariance matches the distribution of the particles passed to it.
The distinction is particularly consequential when the predictive spread
comes mainly from differences between ancestors rather than process noise
within one ancestor.

A scalar Gaussian calculation shows the issue without numerical flow error.
Suppose a map is constructed for prior $\mathcal N(m,P)$ and observation
$y=x+\varepsilon$, $\varepsilon\sim\mathcal N(0,R)$. An exact affine
Gaussian transport is
\begin{equation}
 F_{m,P}(x)=m+\frac{P}{P+R}(y-m)
              +\sqrt{\frac{R}{P+R}}(x-m).
\end{equation}
If its input really has law $\mathcal N(m,P)$, the output has the correct
posterior mean and variance. If the input is instead a particular parent's
transition $X\sim\mathcal N(a,Q)$, direct substitution gives
\begin{align}
 \mathbb E[F_{m,P}(X)]&=m+\frac{P}{P+R}(y-m)
                    +\sqrt{\frac{R}{P+R}}(a-m),\\
 \operatorname{Var}[F_{m,P}(X)]&=\frac{RQ}{P+R}.
\end{align}
The desired conditional posterior has mean
$a+Q(y-a)/(Q+R)$ and variance $QR/(Q+R)$. For $a=m=0$, $Q=1$,
$P=100$, $R=100$ and $y=10$, the map yields mean $5$ and variance $0.5$,
whereas conditioning the actual transition yields approximately $0.0990$
and $0.9901$. The map is correct for its stated Gaussian input; using it as
the exact conditional update for this different input is wrong.

This example identifies a mechanism, rather than measuring how much of a
particular nonlinear run's error it explains. A proposal with this mismatch
can still have a valid importance correction if its actual density is
evaluated. But valid weights cannot guarantee that a finite sample covers
the important region well. Parent-conditioned proposals address the
within-parent mismatch; predictive-observation weighting addresses the
allocation of children between parents. These are distinct repairs.

\subsection{An exact update of an inaccurate incoming mixture}

Even exact conditioning only updates the incoming particle approximation.
If the previous cloud omits a region of the filtering distribution, the
finite-mixture normalizer $N^{-1}\sum_j L_j(y)$ can differ from the true
predictive likelihood. Exact conditional sampling cannot reconstruct
ancestors that are absent. Transport and moment restoration also retain only
specified properties of the cloud; preserving its first two moments does
not determine its tails, multimodality or parameter sensitivity.

The familiar weight diagnostic
\begin{equation}
 \operatorname{ESS}=\frac{1}{\sum_i \bar w_i^2}
\end{equation}
measures concentration of normalized importance weights. With exact
observation-weighted conditional sampling of the finite mixture, the
incremental weights are constant and ESS is $N$. This does not make the
incoming approximation exact, nor does it imply $N$ independent informative
ancestors after earlier resets. Higher ESS has practical value because less
of the represented mass is carried by a few particles. It remains a measure
of weight balance, not of distance to the model posterior or its score.

The SIR controls make this limitation observable. In the pure
conditional-only ablation (the defensive branches are removed), at nominal
$T=50$ the weighted proposal attains ESS $1008$, but its infection-score
RMSE is $95.82$--$100.96$ across the two references. The conditional mixture
with uniform parents gives $31.91$--$32.76$ in the corresponding
three-branch comparison. At noise multiplier $1.25$,
$T=50$, identity-only and conditional-only uniform proposals also have lower
observed noise-score RMSE than the weighted conditional mixture. These
controls prevent a claim of accuracy optimality. They do not remove the
value of the measured reduction in weight concentration.

\subsection{Differentiating the finite calculation is a different test}

The model score is $s(\theta)=\nabla_\theta\log p_\theta(y_{1:T})$.
The implemented analytical calculation differentiates the finite filter
value $\widehat\ell_N(\theta;u)$ for fixed random design $u$, with sampled
categorical labels held locally fixed. A finite-difference comparison away
from selection boundaries checks that branchwise derivative. It does not
test whether the derivative is close to $s(\theta)$, nor whether it equals
the derivative after averaging over parameter-dependent selections.

A two-parent example makes the distinction explicit. Let $U$ be uniform on
$(0,1)$, choose parent 1 if $U<\alpha(\theta)$ and parent 2 otherwise, and
return fixed values $a$ and $b$. For almost every fixed $U$, the selected
value is locally constant in $\theta$, so its branchwise derivative is zero.
Yet
\begin{equation}
 \mathbb E[h_\theta(U)]=\alpha(\theta)a+
               [1-\alpha(\theta)]b,
 \qquad
 \frac{\mathrm d}{\mathrm d\theta}\mathbb E[h_\theta(U)]
       =\alpha'(\theta)(a-b).
\end{equation}
The derivative of the expectation contains the change in selection
probability. This example demonstrates why local derivative parity alone
cannot settle a model-score claim. It does not prove that every observed
error in the full filter is caused by a missing selection contribution:
the full calculation also differentiates densities, normalizations, state
locations and resets, and some terms can cancel under particular identities.
That causal claim requires an isolated comparison.

There is a second distinction even without categorical choices. If a positive
likelihood estimator satisfies $\mathbb E[\widehat Z]=Z$, and differentiation
may pass through the expectation, then
$\mathbb E[\nabla\widehat Z]=\nabla Z$. It does not follow that
\begin{equation}
 \mathbb E\!\left[\frac{\nabla\widehat Z}{\widehat Z}\right]
       =\frac{\nabla Z}{Z}.
\end{equation}
The random denominator is correlated with the numerator. Correctness of an
unnormalized derivative identity therefore does not remove finite-particle
score bias. Nor does a small error in a likelihood value bound its derivative
error. These are reasons to compare likelihood and score separately.

\subsection{What the earlier score repairs did and did not solve}

The observation-sidecar and kernel-mixture investigation exposed a target
distinction before parent conditioning was introduced. An expected
conditional log likelihood is generally different from the log marginal
likelihood. For a fixed mixing measure, Jensen's inequality gives
$\mathbb E[\log g_\theta(y\mid X)]\leq
\log\mathbb E[g_\theta(y\mid X)]$; their derivatives need not agree.
An unbiased gradient identity for the former cannot establish an unbiased
model score for the latter. The research note derives the precise
importance-weighted gradient identity and its normalization requirements.

The subsequent complete kernelized-observation recursion and
all-components resampling reference corrected more of the finite
calculation and passed their declared derivative checks. Nevertheless, the
calibrated resampling construction had higher observed model-score MSE than
the canonical comparator in both tested matrix-Gaussian scopes. That is a
negative result for the tested construction, rather than evidence that the
derivative check failed or that every kernel-mixture proposal must fail.
The later simulation-supervised and neural-guide comparisons at $N=64$
were withdrawn; they cannot support an accuracy ranking or rejection of
guidance. The full derivations and scope-specific evidence are retained in
the research note, rather than treating those discarded runs as support
for the present default.

\subsection{The model-specific failures still requiring explanation}

KSC is a useful counterexample to an ESS-based score argument. Weighted
parents raise mean ESS from about $883$ to $997$ and reduce log-likelihood
RMSE, yet persistence-score RMSE increases from about $0.074$ to $0.407$.
Because the conditional calculation is exact for the implemented
seven-mixture observation target, a generic claim that all remaining error
comes from an unconditioned child is inadequate. Recursive representation,
selection sensitivity and inherited numerical controls remain possible
causes; the completed experiment does not isolate their contributions.

In range--bearing, conditioning is locally Gaussian and the bearing residual
is wrapped. Its conditional proposal is therefore an approximation to a
nonlinear observation update. Weighted parents reduce the observed
process-scale score RMSE relative to the original, but increase the
observation-scale score RMSE; uniform parents have lower RMSE in both
components than weighted parents. Reference Monte Carlo error limits the
precision of this comparison. The experiment does not establish whether
local linearization, allocation, or recursive sensitivity explains the
remaining discrepancy.

Predator--prey is the clearest warning against a universal cross-model
claim: all six weighted-parent score RMSEs exceed the original comparator's
values. The adapted Zhao--Cui local fits are themselves approximate, so
better reference and support checks are needed before assigning the full
gap to the particle proposal. LGSSM supplies the exact-oracle control:
the small, mixed differences show no broad accuracy advantage for either
parent rule. SIR supplies the opposite extreme, where conditioning the
children removes very large observed score errors but additional parent
weighting often trades score accuracy for ESS and likelihood accuracy.

The next discriminating experiment should hold the incoming cloud fixed,
compare exact finite-mixture integration with conditional sampling, and
separate local derivatives from derivatives of expectations over selection.
Only after that distinction is resolved should independently calibrated
recursive runs identify the effect of reset and allocation choices. The
current evidence supports retaining weighted parents as the requested
default and exposing uniform parents as a competitive option. It does not
close model-score correctness, HMC readiness, or accuracy-optimality questions.
